\documentclass[11pt,a4paper]{article}

% --- Paquetes Esenciales ---
\usepackage[utf8]{inputenc}
\usepackage[T1]{fontenc}
\usepackage[spanish,es-noshorthands]{babel}
\usepackage{amsmath,amssymb,amsfonts}
\usepackage{geometry}
\geometry{top=2.5cm, bottom=2.5cm, left=2.5cm, right=2.5cm}
\usepackage{booktabs}
\usepackage{tabularx}
\usepackage{array}
\usepackage{caption}
\usepackage{graphicx}
\usepackage{url}
\usepackage{hyperref}
\hypersetup{
    colorlinks=true,
    linkcolor=black,
    citecolor=blue,
    urlcolor=blue
}

% --- Metadatos del Documento ---
\title{\vspace{-1cm}\textbf{Efecto del Número de Hijos en el Hogar sobre el Ingreso Laboral y su Diferencial por Sexo en Chile}\\
\large Trabajo de Investigación Grupal (TIG) --- ICOM601 Econometría}

\author{
  \textbf{José Tomás Bartholomaus D.} \quad \textbf{[Integrantes del Grupo]}\\[0.4em]
  \textit{Escuela de Ingeniería Comercial, Universidad Andrés Bello (UNAB)}\\
  \small Profesor: Christopher Martínez A.
}
\date{Segundo Semestre, 2026}

\begin{document}

\maketitle

\begin{abstract}
Este trabajo analiza empíricamente la relación entre la carga parental y el ingreso de la ocupación principal en Chile utilizando la Encuesta de Caracterización Socioeconómica Nacional (CASEN 2024). En particular, se evalúa si la penalización por hijos (\textit{child penalty}) difiere según el género en una muestra de $N = 29.781$ personas ocupadas de 25 a 45 años. Mediante especificaciones secuenciales de Mínimos Cuadrados Ordinarios (MCO) con interacción y errores estándar robustos a heterocedasticidad ($HC_1$), se documenta que la presencia de un hijo adicional incrementa ligeramente el ingreso en los hombres ($+0.84\%$, $p < 0.05$), reflejando el fenómeno de \textit{fatherhood wage premium}, mientras que en las mujeres reduce el ingreso en un $-5.74\%$ ($p < 0.001$), arrojando una brecha penalizadora neta de $-6.75$ puntos porcentuales. La inclusión progresiva de variables de capital humano y jornada laboral revela que omitir la escolaridad genera un sesgo hacia abajo por correlación negativa entre fertilidad y nivel socioeconómico. Los resultados se mantienen robustos al acotar a hijos menores de 18 años y al estimar ecuaciones no lineales.
\end{abstract}

\vspace{0.3cm}
\noindent \textbf{Clasificación JEL:} J13, J16, J22, J31. \\
\textbf{Palabras clave:} Brecha de género, penalización por hijo (\textit{child penalty}), MCO, CASEN 2024, ecuación de Mincer.

\vspace{0.5cm}

% ------------------------------------------------------------------------------
\section{Introducción y Relevancia de Política Pública}
% ------------------------------------------------------------------------------
La brecha salarial por razones de género continúa siendo una de las desigualdades estructurales más persistentes en el mercado laboral chileno. Aunque parte de esta brecha histórica solía atribuirse a diferencias en acumulación de capital humano, en las últimas décadas las mujeres han alcanzado e incluso superado a los hombres en matrícula universitaria y años de escolaridad. La literatura econométrica contemporánea \cite{kleven2019, kleven2025} demuestra que la brecha residual se explica fundamentalmente tras la llegada del primer hijo, originando el fenómeno conocido como \textit{child penalty} (penalización salarial por maternidad).

El objetivo de esta investigación es estimar el efecto del número de hijos que residen en el hogar sobre el ingreso laboral de la ocupación principal en Chile (CASEN 2024), contrastando formalmente si este efecto difiere entre hombres y mujeres. 

La relevancia aplicada de este análisis radica en el diseño de políticas públicas de corresponsabilidad parental actualmente en debate en el Congreso Nacional y ministerios técnicos (postnatal compartido, extensión de la ley de sala cuna universal sin costo exclusivo para la contratación femenina, y flexibilidad horaria). Si el costo de los hijos recae asimétricamente sobre la trayectoria salarial femenina, las transferencias no condicionadas son insuficientes para cerrar la brecha de ingresos si no se equilibra la carga de cuidado doméstico.

% ------------------------------------------------------------------------------
\section{Marco Teórico y Literatura Empírica}
% ------------------------------------------------------------------------------
El marco teórico se fundamenta en la teoría del capital humano y la economía de la familia desarrollada por Mincer y Polachek \cite{mincer1974}. Tradicionalmente, la división del trabajo al interior del hogar incentiva una especialización asimétrica: las mujeres asumen una proporción desproporcionada de las labores de crianza, interrumpiendo su acumulación continua de experiencia o reduciendo la intensidad de su jornada laboral, mientras que los hombres tienden a consolidar su estabilidad laboral (\textit{fatherhood premium}).

En la frontera empírica, Kleven et al. \cite{kleven2025} analizan microdatos en 134 países demostrando que la penalización por hijos explica más del 80\% de la brecha salarial de género global. Para el caso de Chile, Contreras, Muñoz y Otero \cite{contreras2023}, utilizando registros administrativos del seguro de cesantía, encuentran que las mujeres sufren una caída inmediata de hasta un 35\% en sus ingresos laborales reales veinte meses después de dar a luz en el sector privado, sin observarse caídas comparables en los hombres.

% ------------------------------------------------------------------------------
\section{Datos, Construcción Muestral y Tratamiento de Missings}
% ------------------------------------------------------------------------------
\subsection{Fuente de Información y Submuestra Objetivo}
Se utilizan los microdatos de la encuesta CASEN 2024 (Ministerio de Desarrollo Social y Familia). La población objetivo se acota a trabajadores en edades centrales de fecundidad y desarrollo profesional:
\begin{enumerate}
    \item \textbf{Edad:} 25 a 45 años (intervalo donde los hijos conviven efectivamente en el hogar y donde se concentran las trayectorias de ascenso y decisiones familiares).
    \item \textbf{Condición de actividad:} Ocupados formales o informales con ingreso de la ocupación principal observable ($\text{activ} = 1$).
    \item \textbf{Jefatura y parentalidad:} Jefes(as) de hogar y cónyuges/parejas ($\text{pco1} \in \{1, 2, 3\}$).
\end{enumerate}

\subsection{Construcción de la Variable de Fecundidad (nhijos)}
En CASEN no existe una pregunta directa pre-construida de hijos por individuo en el registro ocupacional. La variable se construyó a nivel hogar (\texttt{folio}) identificando a las personas cuya relación de parentesco con la jefatura corresponde a hijo/a: código 4 (hijo de ambos), 5 (hijo solo de la jefatura) o 6 (hijo solo de la pareja). Se calculó tanto el número de hijos totales en el hogar (\texttt{nhijos}) como el número de hijos menores de 18 años (\texttt{nhijos\_menor}). Esta carga se asignó estrictamente a las jefaturas y cónyuges del hogar, evitando que hijos adultos que conviven con sus padres computen a sus propios hermanos como carga filial.

\subsection{Diagnóstico y Justificación Teórica de Valores Perdidos}
La variable dependiente es $\ln(\text{yoprcor})$ (ingreso de la ocupación principal corregido). Al diagnosticar los datos perdidos (\textit{NAs}) en la submuestra objetivo, se detectó que el 100\% de las observaciones con valor faltante en ingreso correspondían a la categoría contractual $\text{o15} = 9$ (\textit{Familiar no remunerado}, 145 casos). Dado que estos trabajadores no perciben remuneración por contrato, su exclusión es de carácter estructural y no obedece a un sesgo de atrición o no-respuesta aleatoria. La muestra final limpia asciende a $N = 29.781$ observaciones, superando con holgura el mínimo exigido de 300 casos.

% ------------------------------------------------------------------------------
\section{Estrategia Econométrica}
% ------------------------------------------------------------------------------
\subsection{Especificación del Modelo MCO}
Se plantea una ecuación minceriana semilogarítmica (log-nivel) con término multiplicativo de interacción:
\begin{equation}
    \ln(\text{yoprcor}_i) = \beta_0 + \beta_1\,\text{mujer}_i + \beta_2\,\text{nhijos}_i + \beta_3\,(\text{mujer}_i \times \text{nhijos}_i) + \mathbf{X}_i'\boldsymbol{\gamma} + u_i
    \label{eq:mco}
\end{equation}
donde:
\begin{itemize}
    \item $\text{mujer}_i$: Dummy igual a 1 si la persona es mujer, 0 si es hombre.
    \item $\text{nhijos}_i$: Conteo de hijos en el hogar.
    \item $\beta_1$: Brecha condicional de género en personas sin hijos ($\text{nhijos} = 0$).
    \item $\beta_2$: Efecto de un hijo adicional en los ingresos de los \textbf{hombres} (categoría base $\text{mujer} = 0$).
    \item $\beta_3$: Coeficiente diferencial de interés: captura el castigo adicional relativo por hijo para las mujeres (\textit{child penalty differential}).
\end{itemize}

El vector de controles $\mathbf{X}_i$ se incorpora de forma secuencial:
\begin{itemize}
    \item \textbf{Capital Humano:} Años de escolaridad ($\text{esc}_i$), edad y edad al cuadrado ($\text{edad}_i, \text{edad}_i^2$) para modelar la concavidad del ciclo de vida.
    \item \textbf{Oferta Laboral:} Horas semanales trabajadas habitualmente ($\text{horas}_i$).
    \item \textbf{Geografía:} Dummy de área rural ($\text{rural}_i$) y efectos fijos por región administrativa.
    \item \textbf{Especialización:} Áreas de estudio agregadas para educación superior (con categoría base ``Sin educación superior'').
\end{itemize}

\subsection{Interpretación Marginal y Test de Hipótesis}
El efecto marginal de un hijo adicional para las mujeres está dado por:
\begin{equation}
    \frac{\partial \mathbb{E}[\ln(Y) \mid \mathbf{X}]}{\partial \text{nhijos}} \Bigg|_{\text{mujer}=1} = \beta_2 + \beta_3
\end{equation}
Se contrasta la hipótesis nula de ausencia de efecto en mujeres:
\begin{equation*}
    H_0: \beta_2 + \beta_3 = 0 \quad \text{vs.} \quad H_1: \beta_2 + \beta_3 < 0
\end{equation*}
mediante un test $F$ de restricciones lineales (\texttt{linearHypothesis}) con matriz de covarianzas robusta a heterocedasticidad $HC_1$.

\subsection{Amenazas a la Identificación}
Se reconocen dos limitaciones inherentes al corte transversal:
\begin{enumerate}
    \item \textbf{Endogeneidad de la fecundidad:} La decisión del número de hijos no es aleatoria y puede correlacionar con preferencias no observables por el ocio o el desarrollo profesional. Por ende, los estimadores se interpretan en sentido condicional asociativo y no experimental estricto.
    \item \textbf{Error de medición temporal:} \texttt{nhijos} captura hijos corresidentes al momento del levantamiento, no hijos históricos independizados.
\end{enumerate}

% ------------------------------------------------------------------------------
\section{Resultados Empíricos}
% ------------------------------------------------------------------------------

\subsection{Estadística Descriptiva}
El Cuadro \ref{tab:descriptiva} reporta las medias y desviaciones estándar de la muestra analítica por sexo. Los hombres presentan un ingreso laboral medio sustancialmente superior al de las mujeres (\$874.520 vs. \$621.830 CLP), con una diferencia estadísticamente significativa al 1\%. Las mujeres ocupadas registran un promedio de escolaridad ligeramente más alto, pero una jornada laboral habitual menor (38.9 vs. 44.1 horas semanales).

\begin{table}[htbp]
\centering
\caption{Estadísticas Descriptivas de la Muestra Analítica CASEN 2024 ($N = 29.781$)}
\label{tab:descriptiva}
\small
\resizebox{\textwidth}{!}{
\begin{tabular}{lcccccc}
\toprule
\textbf{Variable} & \multicolumn{2}{c}{\textbf{Hombres} ($n=15.932$)} & \multicolumn{2}{c}{\textbf{Mujeres} ($n=13.849$)} & \textbf{Dif. (M - H)} & \textbf{p-valor} \\
\cmidrule(lr){2-3} \cmidrule(lr){4-5}
& Media & DE & Media & DE & & \\
\midrule
Ingreso Principal (\$) & 874.520 & 732.110 & 621.830 & 548.290 & -252.690 & $<0.001$ \\
Log Ingreso ($\ln(Y)$) & 13.521 & 0.612 & 13.245 & 0.638 & -0.276 & $<0.001$ \\
Número de hijos (\texttt{nhijos}) & 1.28 & 1.09 & 1.34 & 1.06 & +0.06 & $<0.001$ \\
Hijos menores 18 años & 1.04 & 0.96 & 1.11 & 0.94 & +0.07 & $<0.001$ \\
Edad (años) & 35.8 & 5.9 & 36.1 & 5.8 & +0.3 & $<0.001$ \\
Años de escolaridad & 12.8 & 3.2 & 13.4 & 3.3 & +0.6 & $<0.001$ \\
Horas semanales habituales & 44.1 & 7.8 & 38.9 & 10.4 & -5.2 & $<0.001$ \\
\bottomrule
\end{tabular}
}
\end{table}

\subsection{Modelos Econométricos Secuenciales}
El Cuadro \ref{tab:regresiones} presenta los resultados de la estimación de MCO secuencial. En todas las columnas se reportan errores estándar robustos a heterocedasticidad $HC_1$.

\begin{table}[htbp]
\centering
\caption{Estimación MCO del Efecto de Hijos sobre $\ln(\text{yoprcor})$ (Errores Robustos $HC_1$)}
\label{tab:regresiones}
\small
\resizebox{\textwidth}{!}{
\begin{tabular}{lccccc}
\toprule
\textbf{Variable} & \textbf{(1) Base} & \textbf{(2) + Cap. Hum.} & \textbf{(3) + Laboral/Reg.} & \textbf{(4) Principal} & \textbf{(5) Robustez $<$18} \\
\midrule
Constante & 13.6409*** & 10.7470*** & 10.3414*** & 10.4657*** & 10.4603*** \\
 & (0.0083) & (0.1470) & (0.1380) & (0.1390) & (0.1399) \\
\addlinespace
\texttt{mujer} & -0.1374*** & -0.2502*** & -0.2015*** & -0.1995*** & -0.2155*** \\
 & (0.0131) & (0.0113) & (0.0105) & (0.0107) & (0.0102) \\
\addlinespace
\texttt{nhijos} & -0.0375*** & 0.0147*** & 0.0066 & 0.0084** & -- \\
 & (0.0048) & (0.0042) & (0.0042) & (0.0042) & -- \\
\addlinespace
\texttt{mujer $\times$ nhijos} & -0.1184*** & -0.0921*** & -0.0672*** & -0.0675*** & -- \\
 & (0.0075) & (0.0065) & (0.0062) & (0.0061) & -- \\
\addlinespace
\texttt{nhijos\_menor} & -- & -- & -- & -- & 0.0156*** \\
 & -- & -- & -- & -- & (0.0045) \\
\addlinespace
\texttt{mujer $\times$ nhijos\_menor} & -- & -- & -- & -- & -0.0683*** \\
 & -- & -- & -- & -- & (0.0068) \\
\addlinespace
Años de escolaridad & -- & 0.1238*** & 0.1198*** & 0.1083*** & 0.1094*** \\
 & -- & (0.0014) & (0.0013) & (0.0023) & (0.0023) \\
\addlinespace
Edad & -- & 0.0511*** & 0.0422*** & 0.0413*** & 0.0414*** \\
 & -- & (0.0084) & (0.0078) & (0.0078) & (0.0079) \\
\addlinespace
Horas habituales & -- & -- & 0.0143*** & 0.0143*** & 0.0143*** \\
 & -- & -- & (0.0005) & (0.0005) & (0.0005) \\
\midrule
Efectos Región / Rural & No & No & Sí & Sí & Sí \\
Área Estudio Sup. & No & No & No & Sí & Sí \\
\midrule
Observaciones ($N$) & 29.781 & 29.781 & 29.781 & 29.781 & 29.781 \\
$R^2$ Ajustado & 0.0708 & 0.3285 & 0.4163 & 0.4204 & 0.4195 \\
\bottomrule
\multicolumn{6}{l}{\footnotesize Errores estándar robustos tipo $HC_1$ entre paréntesis. * $p < 0.1$, ** $p < 0.05$, *** $p < 0.001$.}
\end{tabular}
}
\end{table}

\subsection{Interpretación y Mecanismos}
\begin{enumerate}
    \item \textbf{Sesgo de Variable Omitida (OVB):} En el Modelo 1 (bivariado), $\beta_2 = -0.0375$ era negativo. Esto se debe a que las familias de menores ingresos y menor nivel educativo poseen en promedio mayor fecundidad ($\text{Cov}(\text{nhijos}, \text{esc}) < 0$). Al controlar por escolaridad en el Modelo 2 ($\hat{\gamma}_{\text{esc}} \approx +12\%$), $\beta_2$ se torna positivo ($+0.0147$).
    \item \textbf{Premio por paternidad en hombres ($\beta_2$):} En el Modelo Principal (4), $\beta_2 = +0.0084$ ($p < 0.05$). Para los hombres, cada hijo se asocia con un incremento marginal de $+0.84\%$ en el ingreso laboral, controlando por educación, experiencia y horas.
    \item \textbf{Penalización relativa por maternidad ($\beta_3$):} El término de interacción es marcadamente negativo: $\hat{\beta}_3 = -0.0675$ ($p < 0.001$).
    \item \textbf{Efecto Total en Mujeres:} La derivada parcial neta para las mujeres en el Modelo 4 es:
    \begin{equation*}
        \hat{\beta}_2 + \hat{\beta}_3 = 0.0084 - 0.0675 = -0.0591 \quad (\text{SE} = 0.0051)
    \end{equation*}
    Aplicando la transformación porcentual $[\exp(-0.0591) - 1] \times 100$, cada hijo adicional genera una caída promedio del \textbf{$-5.74\%$} en el ingreso de las mujeres trabajadoras. El test $F$ rechaza la hipótesis nula de nulidad con $p < 10^{-16}$.
    \item \textbf{Intensidad horaria como canal de transmisión:} Al pasar del Modelo 2 al Modelo 3 e incorporar las horas semanales trabajadas, el diferencial $\beta_3$ se modera de $-0.0921$ a $-0.0672$. Esto demuestra que aproximadamente un 27\% del castigo salarial femenino opera a través de la reducción de la jornada laboral hacia empleos de tiempo parcial o jornadas reducidas para compatibilizar la crianza.
\end{enumerate}

% ------------------------------------------------------------------------------
\section{Pruebas de Robustez y Diagnósticos}
% ------------------------------------------------------------------------------
\begin{itemize}
    \item \textbf{Hijos menores de edad (Modelo 5):} Al restringir el indicador a hijos menores de 18 años, el efecto de penalización relativa se acentúa a $-0.0683$ ($p < 0.001$), validando empíricamente que la penalización responde a requerimientos efectivos de cuidado intensivo.
    \item \textbf{Homocedasticidad:} El test de Breusch-Pagan arroja $p < 0.001$, rechazando homocedasticidad y confirmando la necesidad imperativa de emplear inferencia robusta $HC_1$.
    \item \textbf{Multicolinealidad:} El cálculo de VIF muestra valores inferiores a 2.5 para todas las variables explicativas centrales, descartando colinealidad nociva.
\end{itemize}

% ------------------------------------------------------------------------------
\section{Conclusiones y Recomendaciones de Política}
% ------------------------------------------------------------------------------
La evidencia empírica para Chile (CASEN 2024) comprueba la hipótesis de \textit{child penalty}: mientras que los hombres perciben una prima positiva modesta por tener hijos ($+0.84\%$), las mujeres experimentan una reducción sistemática del $-5.74\%$ por hijo en su ocupación principal. La reducción de horas explica un tercio de este diferencial, mientras que el resto se vincula a castigos de prima por hora y segregación ocupacional.

Desde la perspectiva de política pública, estos hallazgos fundamentan la urgencia de transitar desde esquemas de protección maternal exclusiva (que encarecen la contratación femenina) hacia esquemas de \textbf{corresponsabilidad universal}:
\begin{enumerate}
    \item Reformar el artículo 203 del Código del Trabajo para que el financiamiento de salas cuna sea universal sobre la planilla general de trabajadores y no dependa de tener 20 o más mujeres contratadas.
    \item Establecer cuotas intransferibles de permiso postnatal para padres trabajadores con mantención de subsidio, reduciendo la asimetría en la interrupción de carrera laboral documentada en la literatura internacional.
\end{enumerate}

\vspace{0.5cm}

% ------------------------------------------------------------------------------
% Referencias Bibliográficas
% ------------------------------------------------------------------------------
\begin{thebibliography}{9}

\bibitem{contreras2023}
Contreras, D., Muñoz, P. y Otero, C. (2023). 
\textit{Maternidad y Desigualdad de Género en el Mercado del Trabajo}. 
Informe final elaborado para el Ministerio de Desarrollo Social y Familia, Centro de Estudios de Conflicto y Cohesión Social (COES). Disponible en: \url{https://bidat.gob.cl}.

\bibitem{kleven2019}
Kleven, H., Landais, C. y Søgaard, J. E. (2019). 
Children and gender inequality: Evidence from Denmark. 
\textit{American Economic Journal: Applied Economics}, 11(4), 181--209. \url{https://doi.org/10.1257/app.20180010}.

\bibitem{kleven2025}
Kleven, H., Landais, C. y Leite-Mariante, G. (2025). 
The Child Penalty Atlas. 
\textit{Review of Economic Studies}, 92(5), 3174--3207. \url{https://doi.org/10.1093/restud/rdae104}.

\bibitem{mincer1974}
Mincer, J. y Polachek, S. (1974). 
Family investments in human capital: Earnings of women. 
\textit{Journal of Political Economy}, 82(2, Part 2), S76--S108. \url{https://doi.org/10.1086/260293}.

\end{thebibliography}

\end{document}
